% Transcription — fonds Alexandre Grothendieck, shelfmark 82, batch 4 (pages 61-78).
% Established from the facsimile archives/batches/82/batch-04.pdf, cut from a
% copy of 82.pdf fetched through the project's relay
% (https://grothendieck-relay.onrender.com/source/82.pdf); PDF page k =
% archivists' page 60+k, no cover sheet. Per the skill transcribe-grothendieck.
% The folder is seventy-eight pages in four batches (1-20, 21-40, 41-60,
% 61-78); this is the last, eighteen pages. Batch 3 (pp. 41-60) existed when
% this file was finished and its notation was followed (GP(1) as
% \mathrm{GP}(1), the gothic \mathfrak{X} for the typical object, C_{i_S} for
% the categories, \mathfrak{g}, \widetilde{\mathfrak{g}}, \mathcal{M},
% \mathcal{L}); batches 1-2 did not exist yet.
% Pass: Opus 5.5 (claude-opus-5-5), 2026-09-26 — first pass, unchecked against the pages by a human.
%
% Dating and title. \dating is the folder's own catalogue date, copied
% verbatim; the folder belongs to the group « Géométrie et topologie
% combinatoire » (69-89), but since it has a date of its own the group's range
% is not added. \foldertitle is the catalogue title without its final full stop.
%
% Pages skipped: none. Pages summarised: none. Pages transcribed from his
% hand: 61-78, all in blue felt-tip, fast cursive; p. 78 is half a page and
% ends the folder. Drawings: none.
%
% His pagination: none on the leaves. His numbered runs: the equation run
% begun on p. 60 ((1)-(2), batch 3) continues here from (3) on p. 61 to (61)
% on p. 78. Irregularities left as written: (10) twice on p. 63; (42) twice on
% p. 71 and (43) on both p. 71 and p. 72; the number read (26) on p. 67 is
% written over a 5; several numbers are overwritten (14, 18, 21, 22, 24).
% Roman sections continue those of batch 3 (« I Représentants typiques »,
% « II Opérations de G_0 », pp. 59-60): 5) and 6) of II on pp. 61-62, then
% III (p. 63), IV (p. 65), V (p. 67), VI (p. 69). A « Proposition » on p. 76.
%
% What the batch is about. The « Dictionnaire lié à GP(1) ≃ SO(3) » (cover
% p. 56, batch 3) continued. pp. 61-62: G_0 = GP(1) acting on the typical
% objects — the form q_0 = -det = xy + z^2 on the trace-zero matrices
% sl(2) ≃ E_0 = O^3, on P^1, on (P_0, C_0 = V(q_0)); admission that the maps
% G_0 -> Aut(X_0) are isomorphisms (7). pp. 63-67: the functors of the
% dictionary: III, twisting by a torsor R (8), and C_0 -> C_5 as R -> R/B_0
% with the Borel B_0 (9)-(14); IV, the inverse functors X -> Isom(X_0, X),
% X -> Aut(X), G~ -> G~/Cent, M -> V(M)*/G_m (15)-(20); from forms of GP(1)
% to the others: universal cover, Lie algebras, Bor(G) (21)-(26); V, towards
% forms of SL(2) (27)-(31). pp. 68-78, section VI: a quaternion algebra M
% over S as the key to all the objects — reduced norm and trace, the exact
% sequences (34)-(35') for G, G~, g, g~; the trace form Phi(u,v) = Tr(uv),
% the dualities g~ ≃ g^v, the form q = -det on g~ and its relation to the
% Killing form (45); the basis omega of det g~ and the sign convention,
% discriminant delta(q) = -omega^2 (48)-(49); f : g~ -> g and f' with
% ff' = -2 (51)-(55); P and X = V(q) in terms of M (56); a Proposition on
% Borel subgroups via their Lie algebras (57); and in characteristic 2 the
% morphism X -> P(g~/O_S) ≃ X^(2), the Frobenius, radicial and not a
% monomorphism (58)-(61), where the folder breaks off.
%
% Hardest pages: 68 and 71 (heavily corrected), 73-74 and 77 (connective
% prose largely lost, oblique marginal notes on 74 and 77). Notation fixed
% for the batch: his round L for the unit group V(M)* as \mathcal{L}; his
% round letter for the Lie algebra of a Borel as \mathfrak{L}; double-struck
% V as \mathbb{V}; the affine line of (33) as \mathbb{E}^1.

\documentclass[11pt,a4paper]{article}
% Le préambule commun à toutes les transcriptions du fonds Grothendieck.
%
% Il tient en une idée : **l'incertitude se compose, elle ne se résout pas.**
% Une transcription qui lisse un mot douteux a détruit la seule chose pour
% laquelle on la faisait — un lecteur doit pouvoir savoir, à chaque endroit,
% ce qui était sur le feuillet et ce qui a été ajouté. Les six macros qui
% suivent sont donc l'appareil critique, et non de la décoration.
%
% Les mêmes commandes sont reconnues par `scripts/render.mjs`, qui produit la
% vue de lecture du site. Une macro ajoutée ici doit l'être aussi là-bas,
% faute de quoi le rendu échouera bruyamment — ce qui est voulu : un rendu
% silencieusement faux est pire qu'un rendu absent.

\ProvidesPackage{grothendieck}[2026/08/08 Transcriptions du fonds Grothendieck]

\usepackage{amsmath,amssymb,amsthm}
\usepackage{mathrsfs}
\usepackage{stmaryrd}
% Les diagrammes commutatifs sont partout dans ce fonds : c'est la langue
% même de la géométrie algébrique de Grothendieck, et une transcription qui
% les rendrait en prose ne transcrirait rien.
\usepackage{tikz-cd}
% Français par défaut — c'est la langue des feuillets — mais l'anglais est
% chargé aussi : la traduction et le résumé sont des documents anglais, et la
% typographie française (espaces avant les deux-points) y serait une faute.
% Ces documents basculent par \selectlanguage{english} en tête de corps.
\usepackage[english,french]{babel}
\usepackage{fontspec}
\usepackage{geometry}
\geometry{a4paper,margin=25mm,marginparwidth=28mm}
\usepackage{marginnote}
\usepackage{xcolor}
% ulem, not soul. `soul`'s \st and \ul are built on a character-by-character
% scan that XeLaTeX's font machinery breaks: the document fails with an
% undefined control sequence rather than degrading. ulem does the same two jobs
% and is Unicode-engine safe. `normalem` stops it hijacking \emph, which the
% transcriptions use for Grothendieck's own emphasis.
\usepackage[normalem]{ulem}
\usepackage{hyperref}
\hypersetup{colorlinks=true,linkcolor=black,urlcolor=black,pdfborder={0 0 0}}

% --- Métadonnées du lot -----------------------------------------------------
% Reprises telles quelles par le script de rendu, qui les affiche en tête de
% la vue de lecture. Elles fixent ce que le fichier prétend couvrir : sans
% elles, une transcription flotte sans point d'attache dans un fonds de seize
% mille feuillets.

\newcommand{\folder}[1]{\gdef\grothFolder{#1}}
\newcommand{\batch}[1]{\gdef\grothBatch{#1}}
\newcommand{\leaves}[2]{\gdef\grothFirst{#1}\gdef\grothLast{#2}}
\newcommand{\foldertitle}[1]{\gdef\grothFoldertitle{#1}}
\gdef\grothFolder{?}\gdef\grothBatch{?}\gdef\grothFirst{?}\gdef\grothLast{?}\gdef\grothFoldertitle{}

% --- Le résumé --------------------------------------------------------------
%
% Il ouvre la lecture modernisée, sous le titre « Résumé », et il est composé
% en petit. Ce n'est pas un ornement : c'est le seul passage du document qui
% ne soit pas des mathématiques, et un lecteur doit pouvoir le reconnaître
% avant de l'avoir lu — puis le sauter s'il connaît déjà le sujet, ou s'y
% arrêter s'il ne le connaît pas. Le filet qui le suit marque la reprise du
% corps.

\newenvironment{resume}
  {\small\setlength{\parskip}{0.45em}}
  {\par\medskip\hrule\medskip}

% --- Mots-clés ---------------------------------------------------------------
%
% En anglais, en clôture du résumé de la lecture modernisée : le vocabulaire
% moderne sous lequel chercher ce que les feuillets construisent. Le manifeste
% du site (`npm run manifest`) les extrait pour étiqueter la cote entière —
% cette ligne est la source unique des tags, il n'y a pas de fichier à côté.
\newcommand{\keywords}[1]{\par\smallskip\noindent
  {\footnotesize\textsc{Keywords} --- \textit{#1}}\par}

% --- L'appareil critique ----------------------------------------------------

% Le numéro de feuillet, en marge. C'est l'ancre qui permet de régler un
% désaccord sur une formule en regardant *un* feuillet plutôt qu'une chemise
% de six cents. Le site s'en sert aussi pour tourner les pages du fac-similé
% quand on fait défiler la transcription.
\newcommand{\leaf}[1]{%
  \marginnote{\footnotesize\color{gray!70}\textsf{#1}}%
  \label{leaf:#1}\ignorespaces}

% Illisible : on ne devine pas. Un mot inventé qui se lit comme les autres est
% le pire résultat possible.
\newcommand{\ill}{\textcolor{red!60!black}{[\,\ldots\,]}}

% Lecture douteuse : le mot est donné, et signalé comme incertain.
\newcommand{\uncertain}[1]{\uline{#1}}

% Ajout de l'éditeur — une lettre manquante, un mot sous-entendu.
\newcommand{\add}[1]{\textcolor{blue!50!black}{[#1]}}

% Biffé par Grothendieck lui-même. Ce qu'il a rayé fait partie du document :
% une rature dit souvent où la pensée a changé de direction.
\newcommand{\struck}[1]{\sout{#1}}

% Note du transcripteur, sur l'état matériel du feuillet.
\newcommand{\note}[1]{\par\smallskip{\small\color{gray!60!black}\textit{[#1]}}\par\smallskip}

% Note marginale de Grothendieck, distincte du corps du texte.
\newcommand{\marginal}[1]{\par\smallskip\noindent\hspace{1em}%
  {\small\color{gray!60!black}#1}\par\smallskip}

% --- Titre ------------------------------------------------------------------

\AtBeginDocument{%
  \begin{center}
    {\large\bfseries Fonds Alexandre Grothendieck --- cote n\textsuperscript{o}~\grothFolder}\\[2pt]
    {\itshape\grothFoldertitle}\\[4pt]
    {\small feuillets \grothFirst--\grothLast\ (lot \grothBatch)}\\[2pt]
    {\footnotesize\color{gray!60!black}Transcription établie d'après le fac-similé numérisé
     par l'Université de Montpellier.\\
     Les lectures incertaines sont soulignées, les illisibles notées [\,\ldots\,],
     les ajouts entre crochets.}
  \end{center}
  \vspace{1em}}

\endinput


\folder{82}
\batch{4}
\pages{61}{78}
\dating{[vers 1982]}
\watermark{Édition de démonstration}
\foldertitle{SO(3) ≃ GP(1) : notes manuscrites (s.d.)}

\begin{document}

\page{61}
5) Considérons la forme \struck{déterminant} \add{$\xi \mapsto -\det(\xi)$} sur
$\mathrm{M}_0 = \mathrm{M}_2(\underline{O}_{S_0})$, et considérons sa restriction
au sous-Module $\mathfrak{sl}(2)_{S_0} \subset \mathrm{M}_0$ des matrices de trace
nulle, \struck{le groupe} i.e.\ des matrices
\[
  (3) \qquad \xi = \begin{pmatrix} z & x \\ y & -z \end{pmatrix}
\]
donnant
\struck{$q(\xi) = \det \xi =$}
\[
  (4) \qquad q_0(\xi) \overset{\mathrm{déf}}{=} -\det \xi = xy + z^2
\]
\uncertain{Cette} forme \add{($-\det u$ sur $\mathrm{M}_0 = \mathrm{M}_2(\underline{O}_{S_0})$)}
est évidemment invariante par l'action de $\mathrm{GP}(1)$, déduite par passage
au quotient de celle de $\mathrm{GL}(2)$ -- \ill{}
\[
  (5) \qquad \det(g \xi g^{-1}) = \det \xi
\]
a fortiori sa restriction : $\widetilde{\mathfrak{g}}_0 = \mathfrak{sl}(2)_{S_0}$
l'est. Identifions $\widetilde{\mathfrak{g}}_0$ à $E_0 = \underline{O}^3_{S_0}$ par
\[
  (6) \qquad \left\lbrace
  \begin{array}{l}
    (x, y, z) \longmapsto \begin{pmatrix} z & x \\ y & -z \end{pmatrix} \\[4pt]
    E_0 = \underline{O}^3_{S_0} \xrightarrow{\ \sim\ } \widetilde{\mathfrak{g}}_0 = \mathfrak{sl}(2)_{S_0}
  \end{array}\right.
\]
on trouve une action de $G_0 = \mathrm{GP}(1)$ sur $(E_0, q_0)$. \uncertain{Il} est immédiat
\note{les numéros (3) à (6) poursuivent la suite (1)--(2) ouverte à la page 60, dans « II Opérations de $G_0$ ». « $\xi \mapsto -\det(\xi)$ » est écrit au-dessus de « déterminant », biffé ; la parenthèse « ($-\det u$ \ldots) » est une insertion interlinéaire rattachée à « forme » par un trait. Les deux lettres gothiques sont lues $\mathfrak{sl}(2)$ et $\widetilde{\mathfrak{g}}_0$, comme à la page 60.}

\page{62}
aussi, comme $G_0$ opère trivialement sur $\underline{H}$ \uncertain{module}
inversible (car $\underline{\mathrm{Hom}}\bigl(\mathrm{GP}(1), \mathbb{G}_m\bigr) = 0$)
qu'il opère trivialement sur $\underline{\det}(\widetilde{\mathfrak{g}}_0) \simeq
\det(E_0)$, donc \uncertain{\ill{}} \add{\uncertain{invariance}} \uncertain{en} \uncertain{plus} de $q$.
\note{le début de la page enchaîne sur la dernière phrase de la page 61 (« Il est immédiat \ldots ») ; la lettre soulignée $\underline{H}$ et le mot qui la suit sont peu nets, et l'ordre des mots de la dernière ligne est incertain.}

6) $\mathcal{L}_0 = \mathrm{GL}(2)$ opérant sur $\underline{O}^2_{S_0}$, opère sur
$\mathbb{P}^1 = \mathbb{P}(\underline{O}^2_{S_0})$, et comme $\mathbb{G}_m$ opère
trivialement, on en déduit par passage au quotient une action de
$G_0 = \mathrm{GP}(1)$ \uncertain{grâce} à l'iso (2).

Enfin, l'opération de $G_0$ sur $(E_0, q_0)$ définit une opération sur
$\bigl(P_0 = \check{\mathbb{P}}(E_0),\ C_0 = \mathbb{V}(q_0)\bigr)$.

Nous admettrons \add{(\uncertain{pour ne pas nous éterniser} \ldots)} que tous les
homomorphismes
\[
  (7) \qquad G_0 = \mathrm{GP}(1) \longrightarrow \underline{\mathrm{Aut}}(\mathfrak{X}_0)
\]
sont des \emph{iso}, ce qui nous permet de fonder le dictionnaire que nous voulons
\add{(ou dictionnaire des \uncertain{explicitées} \uncertain{\ill{}})}.
Mais nous allons \struck{\ill{}} expliciter \add{\uncertain{les}} points essentiels.
\marginal{Expliciter le th.\ d'isomorphie}
\note{la note marginale est dans la marge gauche, en face de (7), tenue par un trait vertical. La lettre gothique de $\underline{\mathrm{Aut}}(\mathfrak{X}_0)$ est celle de la page 58. Les deux parenthèses ajoutées sont écrites au-dessus de la ligne et reliées au texte par une accolade.}

\page{63}
\emph{III} Le foncteur allant de $C_{0_S} = \mathrm{Tors}\bigl(\mathrm{GP}(1)_S\bigr)$
vers les huit autres \emph{catégories} est le foncteur « \uncertain{torsion} »
\[
  (8) \qquad R \longmapsto R \overset{\mathrm{GP}(1)_S}{\wedge} \mathfrak{X}_{0_S}
\]
où $\mathfrak{X}_0$ est le représentant typique \add{sur $\mathbf{Z}$} de la
catégorie fibrée envisagée. Le théorème d'isomorphie précédent signifie : dire que
ces foncteurs sont des équiv.\ de cat, \ill{} le \uncertain{groupe} de \ill{}
\uncertain{relatif} aux cat.\ (*).

Dans le cas
\[
  \begin{array}{ccc}
    C_{0_S} & \longrightarrow & C_{5_S} \\
    \shortparallel & & \shortparallel \\
    \mathrm{Tors}\bigl(\mathrm{GP}(1)_S\bigr) & & \mathrm{Drpr}(S)
  \end{array}
\]
on peut expliciter aussi ce foncteur comme
\[
  (9) \qquad R \longmapsto R/B_{0_S}
\]
où
\[
  (10) \qquad B_0 \subset \mathrm{GP}(1) = G_0
\]
est le Borel-type de $\mathrm{GP}(1)$, \struck{stabilisateur du} \add{déduit par}
\add{\uncertain{image en quotient par $\mu_2$ du}} \struck{le point} \struck{\ill{} matrices}
groupe des matrices
\[
  (10) \qquad \widetilde{B}_0 = \left\lbrace \begin{pmatrix} \lambda & \mu \\ 0 & \lambda^{-1} \end{pmatrix}
  \;\middle|\; \lambda \text{ dans } \mathbb{G}_m,\ \mu \text{ dans } \mathbb{G}_a \right\rbrace
\]
\struck{\ill{}} qui est aussi le stabilisateur de la droite $\underline{O}_S e_1$ dans
$\underline{O}^2_S$ -- s'identifie
\note{« III » continue la numérotation romaine de « I Représentants typiques » et « II Opérations de $G_0$ » (pages 59--60). Le numéro (10) est écrit deux fois, pour $B_0$ et pour $\widetilde{B}_0$ ; il est laissé tel quel. Dans (8), le premier terme du second membre est couvert de hachures, et $\mathfrak{X}_{0_S}$ porte la lettre gothique de la page 58. Les deux lignes qui suivent (10) sont très raturées : « déduit par » et la parenthèse sur $\mu_2$ sont écrits au-dessus de mots biffés. Dans (10) bis, une première lettre sous $\lambda$ est surchargée.}

\page{64}
à la section \add{\uncertain{nulle} $s_0$} \struck{$\bullet$} de \struck{$\mathbb{P}$}
$X_0 = \mathbb{P}^1$ (la section $\infty$, correspondant à la droite
$\underline{O}_S e_1$), a comme stabilisateur
\[
  (11) \qquad \widetilde{B}'_0 = \left\lbrace \begin{pmatrix} \lambda & 0 \\ \mu & \lambda^{-1} \end{pmatrix}
  \;\middle|\; \lambda \text{ dans } \mathbb{G}_m,\ \mu \text{ dans } \mathbb{G}_a \right\rbrace
\]
Comme l'action de $G_0$ sur $X_0$ \struck{définit} \add{\uncertain{et}} \uncertain{la}
section $s_0$ de $X_0$ définit
\[
  (12) \qquad \begin{array}{c} g \longmapsto g \cdot s_0 \\ G_0 \longrightarrow X_0 \end{array}
\]
qui par passage au quotient définit un \emph{iso}
\[
  (13) \qquad G_0/B_0 \longrightarrow X_0 ,
\]
\struck{on trouve} compatible avec les op.\ de $G_0$, opérant par translations
\uncertain{i.e.}\ \uncertain{sur le} \uncertain{premier membre}, on trouve, en prenant l'image
inverse sur $S$ et en tordant l'iso obtenu par le torseur $R$, l'iso.
\[
  (14) \qquad R \wedge X_{0_S} \simeq R/B_{0_S}
\]
qui établit l'interprétation (9).

Le foncteur allant \add{de chacune des huit} des catégories $C_{1_S} \ldots$ vers
$\mathrm{Tors}\bigl(\mathrm{GP}(1)_S\bigr)$ est le foncteur \uncertain{composé} de
\note{le numéro (14) est surchargé. Le stabilisateur (11) est celui de la section $\infty$, alors que la phrase qui l'introduit nomme d'abord la section nulle $s_0$ : c'est ainsi sur la page. « de chacune des huit » est écrit au-dessus de la ligne et renvoie à « catégories ».}

\page{65}
\[
  (15) \qquad \mathfrak{X} \longmapsto \underline{\mathrm{Isom}}(\mathfrak{X}_0, \mathfrak{X})
\]
vers \struck{les} $\mathrm{Tors}\bigl(\underline{\mathrm{Aut}}(\mathfrak{X}_{0_S})\bigr)$,
suivi de l'équivalence
\[
  (16) \qquad \mathrm{Tors}\bigl(\underline{\mathrm{Aut}}(\mathfrak{X}_{0_S})\bigr) \simeq
  \mathrm{Tors}\bigl(\mathrm{GP}(1)_S\bigr)
\]
déduite de l'iso \struck{\ill{}}
\[
  (17) \qquad \mathrm{GP}(1)_S \xrightarrow{\ \sim\ } \underline{\mathrm{Aut}}(\mathfrak{X}_{0_S})
  \ \bigl(\simeq \underline{\mathrm{Aut}}(\mathfrak{X}_0)_S\bigr)
\]
déduit de (7) par chgt de base.
\uncertain{Également} \uncertain{forme} \uncertain{particulière} du foncteur « \uncertain{tordu} »
dans le cas de $C_{0_S} \to C_{1_S}$, cf ci-dessous.
\note{dans (15), le premier $\mathfrak{X}$ porte un indice $0$ qui paraît de trop ; il est gardé comme sur la page.}

\emph{IV} Le foncteur allant de la catégorie \struck{$C_{1_S}$ des} \struck{formes de
$\mathrm{GP}(1)_S$ vers} des \struck{catégories} $C_{0_S}, C_{2_S}$ etc vers la
catégorie $C_{1_S}$ des formes de $\mathrm{GP}(1)$ s'explicite par
\[
  (18) \qquad \mathfrak{X} \longmapsto \underline{\mathrm{Aut}}(\mathfrak{X})
\]
\struck{Pour expliciter les foncteurs}
Dans le cas \struck{des} $C'_{1_S}$ des formes $\widetilde{G}$ de $\mathrm{SL}(2)$,
il y a aussi une description plus terre à terre :
\[
  (19) \qquad \widetilde{G} \longmapsto \widetilde{G}/\mathrm{Cent}(\widetilde{G})
\]
(\emph{NB} le centre de $\widetilde{G}$ est toujours \struck{\ill{}} \uncertain{isom.} à
$\mu_2$ (de façon unique car $\underline{\mathrm{Aut}}(\mu_2) = 1$)).
\add{Pour le cas de $\underline{\mathrm{Quat}}(\underline{O}_S) \to \mathrm{Fo}(G) \ldots$}
\[
  (20) \qquad \mathcal{M} \longmapsto \check{\mathbb{V}}(\mathcal{M})^{*}/\mathbb{G}_{m_S}
\]
\note{« IV » est précédé d'un trait horizontal qui barre la largeur de la page, séparant le paragraphe III du paragraphe IV. Le numéro (18) est surchargé. La ligne « Pour le cas de \ldots » est écrite en petit, entre les lignes, au-dessus de (20).}

\page{66}
Pour les foncteurs en sens inverse, allant de $C_{1_S} = \mathrm{Fo}\bigl(\mathrm{GP}(1)_S\bigr)$
vers les autres catégories, on peut, soit prendre par le torseur
$R = \underline{\mathrm{Isom}}(G_{0_S}, G)$ pour tordre $\mathfrak{X}_{0_S}$
\[
  (21) \qquad G \longmapsto \underline{\mathrm{Isom}}(G_{0_S}, G) \overset{G_{0_S}}{\wedge} \mathfrak{X}_{0_S} ,
\]
soit trouver, dans divers cas, des descriptions plus directes. Par exemple
\[
  (22) \qquad G \longmapsto \widetilde{G} = \text{Rev.\ simplement connexe de } G
\]
défini comme le schéma en groupes lisse à fibres \struck{\ill{}} \add{\uncertain{géom.\ connexes}} simplement connexes, muni d'une isog.\ $\widetilde{G} \to G$ -- lequel est défini de façon \emph{unique} \ldots

Le foncteur vers les Algèbres de Lie est
\[
  (23) \qquad \begin{array}{c}
    G \longmapsto \mathfrak{g} = \underline{\mathrm{Lie}}(G) \\
    \mathrm{Fo}\bigl(\mathrm{GP}(1)_S\bigr) \xrightarrow{\ \sim\ } \mathrm{Fo}\bigl(\mathfrak{gp}(1)_S\bigr)
  \end{array}
\]
\[
  (24) \qquad \begin{array}{c}
    G \longmapsto \widetilde{\mathfrak{g}} = \underline{\mathrm{Lie}}(\widetilde{G}) \\
    \mathrm{Fo}\bigl(\mathrm{GP}(1)_S\bigr) \xrightarrow{\ \sim\ } \mathrm{Fo}\bigl(\mathfrak{sl}(2)_S\bigr)
  \end{array}
\]
Le foncteur
\[
  G \longmapsto (E, q, \omega)
\]
peut se factoriser par $G \mapsto \widetilde{\mathfrak{g}}$, cf plus bas.
\note{les numéros (21), (22) et (24) sont surchargés. La glose sur le revêtement simplement connexe est écrite à droite, sous (22), en lignes courtes ; « géom.\ connexes » y est écrit au-dessus d'un mot biffé. Les « $\sim$ » au-dessus des flèches de (23) et (24) sont surchargés.}

\page{67}
Le foncteur
\[
  G \longmapsto X
\]
s'explicite comme
\[
  (25) \qquad G \longmapsto \underline{\mathrm{Bor}}(G)
\]
associant à $G$ le \struck{\ill{}} schéma des ses sous-\emph{groupes} de Borel $B$
i.e.\ \struck{\ill{}} des ss-groupes tels que \add{$(G, B)$} \struck{$(B, G)$} soit
loc.\ isom.\ \struck{\ill{}} à $(G_{0_S}, B_{0_S})$. On obtient un \uncertain{isom.},
si $G$ et $X$ se correspondent -- donc $G \simeq \underline{\mathrm{Aut}}(X)$ -- un
\uncertain{isom.}\ can.
\[
  (26) \qquad X \xrightarrow{\ \sim\ } \underline{\mathrm{Bor}}(G)
\]
en associant à une section \add{loc.} de $X$ son stabilisateur dans $G$.
\note{le numéro (25) est surchargé ; celui de la seconde formule, lu (26), est écrit sur un 5. Devant $(G_{0_S}, B_{0_S})$, un premier essai, qui semble commencer par une parenthèse et un $G_0$, est couvert de hachures sur deux lignes.}

\emph{V} Les foncteurs vers les formes de $\mathrm{SL}(2)_S$ s'\struck{f}obtiennent
soit en factorisant par $\mathrm{Tors}\bigl(\mathrm{GP}(1)_S\bigr)$
\[
  (27) \qquad \mathfrak{X} \longmapsto R = \underline{\mathrm{Isom}}(\mathfrak{X}_{0_S}, \mathfrak{X})
  \longmapsto R \overset{\mathrm{GP}(1)_S}{\wedge} \widetilde{G}_{0_S} ,
\]
soit plus \uncertain{naturellement} \ill{}
\[
  (28) \qquad \mathfrak{X} \longmapsto G = \underline{\mathrm{Aut}}_S(\mathfrak{X})
  \longmapsto \widetilde{G} = \text{Rev.\ simpl.\ conn.\ } G
\]
\note{dans (27), le dernier facteur est écrit sur une lettre raturée et surmonté d'un tilde ; il est lu $\widetilde{G}_{0_S}$.}

\page{68}
\uncertain{Dans le cas de} \struck{$\mathrm{Quat}(S)$}
\struck{\ill{}}
\[
  \mathrm{Quat}(S) \longrightarrow \mathrm{Fo}\bigl(\mathrm{SL}(2)_S\bigr)
\]
i.e.\ \uncertain{donné par}
\[
  (29) \qquad \mathcal{M} \longmapsto \widetilde{G} \simeq
  \mathrm{Ker}\Bigl(\check{\mathbb{V}}(\mathcal{M})^{*} \xrightarrow{\ \det\ } \mathbb{G}_m\Bigr)
\]
\note{la première ligne est en partie biffée, et la ligne qui la suit, illisible, l'est entièrement. Dans (29), un premier essai après « $\simeq$ » est couvert de hachures ; sous la flèche « $\det$ », une remarque soulignée : « $=$ norme réduite ».}

Parmi les foncteurs de $\mathrm{Fo}\bigl(\mathrm{SL}(2)_S\bigr)$ vers les autres
catégories, \ill{} \uncertain{sont} \uncertain{pour} $\mathrm{Tors}\bigl(\mathrm{GP}(1)\bigr)$
\ill{} vers $\mathrm{Fo}\bigl(\mathrm{GP}(1)_S\bigr)$, \uncertain{suivis} des
\uncertain{certains}, \uncertain{en} \uncertain{spéciales} \ill{}, p.ex.
\[
  \mathrm{Fo}\bigl(\mathrm{SL}(2)_S\bigr) \longrightarrow \mathrm{Fo}\bigl(\mathfrak{sl}(2)_S\bigr)
\]
\[
  (30) \qquad \widetilde{G} \longmapsto \widetilde{\mathfrak{g}} = \underline{\mathrm{Lie}}(\widetilde{G})
\]
\[
  (31) \qquad \widetilde{G} \longmapsto X \simeq \underline{\mathrm{Bor}}(\widetilde{G})
\]

Pour la relation \struck{\ill{}} entre $G, \widetilde{G}, \mathfrak{g},
\widetilde{\mathfrak{g}}, (E, q, \omega)$, c'est $\mathcal{M}$ qui donne
\uncertain{peut-être} la meilleure clef \uncertain{pour les} \uncertain{décrire} tous
-- ainsi que $X$ et $(P, C)$, \uncertain{tant qu'à} faire ! -- et pour comprendre les
relations entre tous ces objets. On verra ainsi que $\mathfrak{g}$ et
$\widetilde{\mathfrak{g}}$ sont \ill{} \uncertain{duaux} l'un de l'autre, par
$E \simeq \widetilde{\mathfrak{g}}$ canoniquement, i.e.\ quand \ill{}
\note{dans (31), l'argument de $\underline{\mathrm{Bor}}$ est surchargé ; après « relation », un passage biffé se termine par « $\mathcal{M}$, $q$ » couvert de hachures, et « Pour » est écrit sur un premier mot.}

\page{69}
\uncertain{par} $\widetilde{\mathfrak{g}}$, on \uncertain{verra} \uncertain{des} \uncertain{interprétations}
remarquables de $\widetilde{\mathfrak{g}} \to \mathfrak{g}$ et $\widetilde{G} \to G$ etc.\ \ldots

\emph{VI} Soit donc $\mathcal{M}$ une Algèbre de quaternions sur $S$ -- i.e.\ une
Alg.\ loc.\ \add{(fppf)} iso.\ à $\mathrm{M}_2(\underline{O}_S)$. On sait que
\[
  (32) \qquad \left\lbrace
  \begin{array}{l}
    \text{Centre de } \mathcal{M} = \underline{O}_S \cdot 1 \\
    \bigl(\text{Centre de } \mathcal{M}^{*} = \underline{O}_S^{*} \cdot 1 \text{ ou encore }
    \bigl(\text{Centre de } \mathcal{L} = \mathbb{V}(\mathcal{M})^{*}\bigr) = \mathbb{G}_m \subset \mathcal{L}(\mathcal{M})\bigr)
  \end{array}\right.
\]
\note{dans les deux lignes de (32), le premier « $\mathcal{M}$ » est suivi de lettres couvertes de hachures ; la seconde ligne s'ouvre sur une double parenthèse.}

\struck{\ill{}} \uncertain{Puis} on a des fonctions « \ill{} réduite » et
« trace réduite »
\[
  (33) \qquad \mathbb{V}(\mathcal{M}) \overset{\det}{\underset{\mathrm{tr}}{\rightrightarrows}} \mathbb{E}^{1}
\]
définies par des conditions évidentes (par \struck{\ill{}} \add{localisation} :
fonctions $\det$ et $\mathrm{tr}$ sur $\mathrm{M}_2(\underline{O}_S)$, \struck{\ill{}}
qui sont invariantes par automorphismes)
\note{à gauche, un premier numéro (peut-être « (33) ») est couvert de hachures, suivi d'un trait. Le but de (33) est écrit sur un $\underline{O}_S$ surchargé et se lit $\mathbb{E}^1$ (la droite affine) ; « localisation » remplace un mot biffé.}

\marginal{\emph{NB} $\det(uv) = \det u \det v$, $\left\lbrace\begin{array}{l}\mathrm{tr} \text{ est } \underline{O}_S\text{-linéaire} \\ \det \text{ est quadratique}\end{array}\right.$ ; $\det(u + \lambda 1) =$ \struck{\ill{}} $\lambda^2 + \lambda \mathrm{Tr}\, u + \det u$. L'appl.\ tangente de $\det$ en la section $1$ est $\mathrm{tr}$}
\note{cet encadré est à droite de la page, tenu par un trait vertical, en face des lignes qui précèdent (33).}

\uncertain{On} \uncertain{obtient} des suites exactes
\[
  (34) \qquad 1 \longrightarrow \mathbb{G}_m \xrightarrow{\text{incl.\ du centre}}
  \mathcal{L}(\mathcal{M}) \longrightarrow G \longrightarrow 1
\]
\[
  (35) \qquad 1 \longrightarrow \widetilde{G} \longrightarrow \mathcal{L}(\mathcal{M})
  \xrightarrow{\ \det\ } \mathbb{G}_m \longrightarrow 1
\]
et sur les algèbres de Lie
\[
  (34') \qquad 0 \longrightarrow \underline{O}_S \xrightarrow{\lambda \mapsto \lambda \cdot 1}
  \mathcal{M} \longrightarrow \mathfrak{g} \longrightarrow 0
\]
\[
  (35') \qquad 0 \longrightarrow \widetilde{\mathfrak{g}} \longrightarrow \mathcal{M}
  \xrightarrow{\ \mathrm{tr}\ } \underline{O}_S \longrightarrow 0 .
\]
\note{le numéro « VI » est souligné d'un trait au-dessus ; c'est le seul numéro romain de ce lot, et le titre qu'il ouvre n'est pas écrit. $\mathcal{L}(\mathcal{M})$, lettre ronde, est le groupe $\mathbb{V}(\mathcal{M})^{*}$ des unités de $\mathcal{M}$, comme dans (32).}

\page{70}
Ces suites exactes déterminent $G, \widetilde{G}, \mathfrak{g}, \widetilde{\mathfrak{g}}$
chacun directement en termes de $\mathcal{L}(\mathcal{M})$, et $\det$ (pour les
groupes) et $\mathcal{M}$ et $\mathrm{tr}$ (pour les algèbres de Lie).
\uncertain{D'un} \uncertain{autre côté}, $\mathcal{M}$ \uncertain{est} considéré comme
Algèbre de Lie \uncertain{grâce à} la forme trace.

Les hom.\ canoniques
\[
  (36) \qquad \left\lbrace
  \begin{array}{l}
    \widetilde{G} \longrightarrow G \\
    \widetilde{\mathfrak{g}} \xrightarrow{\ f\ } \mathfrak{g}
  \end{array}\right.
\]
sont les composés
\[
  (37) \qquad \widetilde{G} \xhookrightarrow{\ \mathrm{incl.}\ } \mathcal{L}(\mathcal{M})
  \xrightarrow{\ \mathrm{proj}\ } G
\]
\[
  (38) \qquad \widetilde{\mathfrak{g}} \xhookrightarrow{\ \mathrm{incl}\ } \mathcal{M}
  \xrightarrow{\ \mathrm{proj}\ } \mathfrak{g}
\]

Considérons sur $\mathcal{M}$ la forme \struck{quadratique} \add{bilinéaire}
\struck{\ill{}}.
\[
  (39) \qquad \Phi(u, v) = \mathrm{Tr}(uv)
\]
On voit alors que $\varphi$ est symétrique et \uncertain{non} \emph{dégénérée} sur
chaque fibre, permettant donc d'identifier $\mathcal{M}$ à son dual
\uncertain{i.e.}
\[
  (40) \qquad \mathcal{M} \simeq \check{\mathcal{M}}
\]
L'orthogonal de $\underline{O}_S$ est $\widetilde{\mathfrak{g}}$, et on trouve ainsi
un accouplement \uncertain{induisant}
\note{la lettre de (39) est surchargée ; la prose qui suit l'écrit $\varphi$. « bilinéaire » est écrit au-dessus de « quadratique », biffé.}

\page{71}
\[
  (41) \qquad \widetilde{\mathfrak{g}} \times \mathfrak{g} \longrightarrow \underline{O}_S
\]
dont on déduit les isom.\ de Modules :
\[
  (42) \qquad \widetilde{\mathfrak{g}} \simeq \check{\mathfrak{g}}, \qquad
  \mathfrak{g} \simeq \check{\widetilde{\mathfrak{g}}}
\]
\struck{\ill{}} La forme quadratique \add{$Q(u) = -\det u$} sur $\mathcal{M}$ est
lisse, et la forme bilinéaire associée est
\[
  (43) \qquad \left\lbrace
  \begin{array}{l}
    \Phi_1(u, v) = -\mathrm{Tr}\, u \,\mathrm{Tr}\, v + \mathrm{Tr}(uv) = -\mathrm{Tr}\, u \,\mathrm{Tr}\, v + \Phi(u, v) \\
    \text{donc } -\det(u + v) + \det u + \det v = \Phi_1(u, v) \\
    \qquad = -\mathrm{Tr}\, u \,\mathrm{Tr}\, v + \mathrm{Tr}(uv)
  \end{array}\right.
\]
qui est donc \emph{également} dualisante. Sur $\widetilde{\mathfrak{g}}$ elle
coïncide avec $\varphi$, et comme $\mathrm{Tr}(1) = 2$, on a
\[
  \Phi_1(1, u) = -\Phi(1, u) = -\mathrm{Tr}\, u
\]
donc l'orth.\ de $\underline{O}_S \cdot 1$ pour $\varphi_1$ est encore
$\widetilde{\mathfrak{g}}$. L'accouplement \add{$\mathfrak{g} \times \widetilde{\mathfrak{g}} \to \underline{O}_S$}
\struck{\ill{}} déduit de $\varphi_1$ par passage au quotient \ill{}
\struck{\ill{}} \add{\uncertain{est égal au signe près}} à (41) défini par $\Phi$.
\struck{\ill{}} \add{D'ailleurs} $\Phi$ et $\Phi_1$ ont même restriction à
$\widetilde{\mathfrak{g}}$, \struck{\ill{}} soit $\varphi$
\[
  (42) \qquad \left\lbrace
  \begin{array}{l}
    \varphi = \Phi \,|\, \widetilde{\mathfrak{g}} \times \widetilde{\mathfrak{g}}
    = \Phi_1 \,|\, \widetilde{\mathfrak{g}} \times \widetilde{\mathfrak{g}} :
    \widetilde{\mathfrak{g}} \times \widetilde{\mathfrak{g}} \to \underline{O}_S \\
    \varphi(u, v) = \mathrm{Tr}\, uv
  \end{array}\right.
\]
\note{le numéro (42) sert deux fois sur cette page, et (43) resservira à la page 72 : c'est sa numérotation, laissée telle quelle. Dans (42), un petit signe de correction surmonte le premier $\check{\mathfrak{g}}$. La fin du paragraphe est très raturée : les insertions « est égal au signe près » et « D'ailleurs » sont écrites au-dessus de mots biffés, et l'ordre de lecture de la ligne est incertain.}

\page{72}
Considérons la forme quadratique $q$ sur $\widetilde{\mathfrak{g}}$
\[
  (43) \qquad \left\lbrace
  \begin{array}{l}
    q = Q \,|\, \widetilde{\mathfrak{g}} : \widetilde{\mathfrak{g}} \longrightarrow \underline{O}_S \\
    q(u) = -\det u
  \end{array}\right.
\]
alors $\varphi$ est donc associée à la forme $q$
\[
  (44) \qquad \left\lbrace
  \begin{array}{l}
    q(x + y) = q(x) + q(y) + \varphi(x, y) \\
    \varphi(x, x) = 2 q(x)
  \end{array}\right.
\]
d'ailleurs $q$ est lisse -- $\varphi$ est dualisante sauf en car.~2.

[La relation avec la forme de Killing est
\[
  (45) \qquad \left\lbrace
  \begin{array}{l}
    \mathrm{Tr}\bigl(\mathrm{ad}_{\widetilde{\mathfrak{g}}}(u)\, \mathrm{ad}_{\widetilde{\mathfrak{g}}}(v)\bigr) = 4 \varphi(u, v) \\
    \mathrm{Tr}\bigl(\mathrm{ad}_{\widetilde{\mathfrak{g}}}(u)^2\bigr) = 8 q(u)
  \end{array}\right.
\]
pour $u, v$ sections de $\widetilde{\mathfrak{g}}$.]

On définit aussi une base
\[
  (46) \qquad \omega \in \Gamma\, \underline{\det}\, \widetilde{\mathfrak{g}} \simeq \Gamma\, \underline{\det}\, \mathcal{M}
\]
(le dernier iso provenant de la suite exacte (35') \ill{} \ill{}) compte tenu de
l'isom.\ \uncertain{local}
\[
  \mathcal{M} \simeq \underline{\mathrm{End}}(V) \simeq \check{V} \otimes V
\]
et de l'isomorphisme canonique
\[
  (47) \qquad \det(\check{V} \otimes V) \simeq \det \check{V} \otimes \det V
  \simeq (\det V)^{-1} \otimes \det V \simeq \underline{O}_S
\]
\note{les crochets autour de (45) sont les siens. Dans (47), le troisième terme est souligné sur la page.}

\page{73}
qui est évidemment stable par $\mathrm{GL}(V)$ donc par $\mathrm{GP}(V)$ -- donc on
déduit une base $\omega$ de $\underline{\det}\, \mathcal{M}$, donc de
$\det \widetilde{\mathfrak{g}}$. Il est vrai que dans l'iso (46) il y a une
convention de \emph{signe} à faire -- à défaut, il y a l'ambiguïté du \emph{signe}
(qui pourrait définir une \struck{\ill{}} \ill{} canonique \ill{} \ill{} -- cette
ambiguïté était évidente, car \uncertain{sur} $\mathbf{Z}$ \ill{} inversibles
\ill{} exactement deux bases, \uncertain{différant} \uncertain{par} le \emph{signe} --
et l'invariance par \struck{$\mathrm{GP}(1)$} $G_0$ est claire.
Mais $\omega^{\otimes 2}$ ne dépend plus de cette convention, ce qui dispense de
lever cette ambiguïté :
\[
  (48) \qquad \underbrace{\delta'(q)}_{\text{discriminant divisé}} = -\,\omega^{\otimes 2}
\]
d'où
\[
  (49) \qquad \delta(\varphi) = \delta(q) = -2\,\omega^{\otimes 2}
\]
\note{la parenthèse ouverte après « \emph{signe} » n'est pas refermée. Sous $\delta'(q)$, souligné, il écrit « discriminant » puis « divisé » ; les \emph{signe} sont soulignés sur la page.}

\page{74}
On \struck{posera}
\[
  (50) \qquad (E, q, \omega) = (\widetilde{\mathfrak{g}}, q, \omega) \qquad \text{i.e. } E = \widetilde{\mathfrak{g}}
\]
\struck{On peut aussi définir,}

\emph{NB} L'hom.\ \struck{fondamental}
\[
  (51) \qquad f : \widetilde{\mathfrak{g}} \longrightarrow \mathfrak{g}
\]
de (36) \uncertain{n'est} autre que l'hom.\ $\widetilde{\mathfrak{g}} \to
\check{\widetilde{\mathfrak{g}}}$ associé à la forme bilinéaire $\varphi$,
compte tenu de l'iso $\check{\widetilde{\mathfrak{g}}} \simeq \mathfrak{g}$.
Ceci signifie en effet que
\[
  \underbrace{\langle \bar{u}, f(v) \rangle}_{\mathrm{Tr}\, uv}
  = \underbrace{\varphi(u, v)}_{\mathrm{Tr}\, uv}
  \qquad u, v \in \Gamma(S, \widetilde{\mathfrak{g}})
\]
Ceci dit, les lois de crochet sur $\widetilde{\mathfrak{g}}$ \uncertain{et} \ill{}
\ill{} $\mathfrak{g}$ sont alors déductibles canoniquement \uncertain{par ailleurs} des
couples $(\varphi, \omega)$ d'une forme bil.\ sym.\ \uncertain{dans}
$\widetilde{\mathfrak{g}}$ et d'une base $\omega$ de $\det \widetilde{\mathfrak{g}}$
(c'est ici que la convention de \emph{signe} \uncertain{s'avère} \uncertain{importante}
-- on prendra
\[
  (52) \qquad \boxed{\ \omega = \widetilde{X} \wedge \widetilde{H} \wedge \widetilde{Y}
  = -\,\underset{e_1}{\widetilde{X}} \wedge \underset{e_2}{\widetilde{Y}} \wedge \underset{e_3}{\widetilde{H}}\ }
\]
\ill{}
\marginal{Sans garantie -- il \struck{\ill{}} \uncertain{faut} vérifier \uncertain{le texte} \ldots}
\note{« On posera » : « posera » est biffé et le verbe n'est pas remplacé. Dans la formule qui suit (51), les deux accolades sont des guillemets de répétition (« \,"\, ») posés sous $\langle \bar{u}, f(v)\rangle$ et $\varphi(u,v)$, chacun au-dessus de « $\mathrm{Tr}\, uv$ ». La formule (52) est encadrée ; sous $\widetilde{X}, \widetilde{Y}, \widetilde{H}$ du second membre, des guillemets de répétition renvoient à $e_1, e_2, e_3$, base de $E_0$ à la page 60. La note « Sans garantie \ldots » est dans la marge gauche, en face de (52), et le mot qui suit la formule (peut-être « On ») est le début d'une phrase qui ne se poursuit pas.}

\page{75}
\struck{\ill{}} \add{On} aussi \uncertain{explicitera} dans notre contexte actuel l'hom.
\[
  (53) \qquad f' : \mathfrak{g} \longrightarrow \widetilde{\mathfrak{g}}
\]
(s'identifiant à ${\textstyle\bigwedge^2} f$) tel que
\[
  (54) \qquad \left\lbrace
  \begin{array}{l}
    f f' = -2\, \mathrm{id}_{\mathfrak{g}} \\
    f' f = -2\, \mathrm{id}_{\widetilde{\mathfrak{g}}}
  \end{array}\right.
\]
en posant
\[
  (55) \qquad f'(\bar{u}) = (\mathrm{Tr}\, u)\, 1 - 2u
\]
où $u$ est une section loc.\ de $\mathcal{M}$, et $\bar{u}$ son image par
$\mathcal{M} \to \mathfrak{g}$ ; comme $\mathrm{Tr}\, 1 = 2$ cette application
\struck{\ill{}} $\mathcal{M} \to \mathcal{M}$ est bien à valeurs dans
$\widetilde{\mathfrak{g}}$, et nulle sur $\underline{O}_S \cdot 1$ donc se factorise
par $\mathfrak{g}$, \uncertain{enfin} si $u$ est dans $\widetilde{\mathfrak{g}}$ on
trouve $f'f(u) = -2u$, enfin \ill{}, $f f'(\bar{u}) =$ \struck{\ill{}}
$\overline{f'(u)} = -2\bar{u}$ ok.

\uncertain{Il} reste à voir comment s'explicitent en termes de $\mathcal{M}$ les
objets $X$, et $P \supset X$. Bien sûr, on aura
\[
  (56) \qquad \left\lbrace
  \begin{array}{l}
    P = \check{\mathbb{P}}(\widetilde{\mathfrak{g}}) \simeq \mathbb{P}(\mathfrak{g}) \\
    X = \mathbb{V}(q) \subset P
  \end{array}\right.
\]
\note{dans (56), l'argument de $\check{\mathbb{P}}$ est surchargé (un premier $\mathfrak{g}$ changé en $\widetilde{\mathfrak{g}}$, semble-t-il). Dans la marge gauche, en face de (56), une note commençant par « NB » est entièrement biffée et illisible.}

\page{76}
La relation \struck{\ill{}} entre cette description du fibré en droites
projectives $X$, et la description comme $\underline{\mathrm{Bor}}(G) \simeq
\underline{\mathrm{Bor}}(\widetilde{G})$, est explicitée dans la

\emph{Proposition.} Les ss-\emph{groupes} de Borel \add{$B$} de $G$ correspondent
biunivoquement par image inverse \add{$\widetilde{B}$} : \struck{ceux}
\add{$B_1$} de $\mathcal{L} = \mathbb{V}(\mathcal{M})^{*}$ \ill{} \add{$\widetilde{B}$}
de $\widetilde{G}$. Un ss-\emph{groupe} de Borel est connu quand on connaît son
algèbre de Lie $\mathfrak{L}$ (qui est un ss-fibré vectoriel de corang~1 de
$\mathfrak{g}$), \uncertain{ou} celle $\mathfrak{L}_1$ de $\mathcal{L}$ qui lui
correspond (qui est un ss-fibré \struck{de rang~2} \add{vectoriel} de corang~1)
image inverse de $\mathfrak{L}$, dans $\mathcal{M} = \underline{\mathrm{Lie}}(\mathcal{L})$,
\struck{ou celle} ou \emph{en car.\ $\neq 2$} quand on connaît celle
$\widetilde{\mathfrak{L}}$ de $\widetilde{B}$ dans $\widetilde{\mathfrak{g}}$, qui
n'est autre que $\mathfrak{L}_1 \cap \widetilde{\mathfrak{g}}$ -- (et \ill{} \ill{}
\ill{}, \uncertain{puisqu'alors} $f : \widetilde{\mathfrak{g}} \xrightarrow{\sim}
\mathfrak{g}$ !). \struck{$\widetilde{B}$} \add{En termes} de $\widetilde{\mathfrak{L}}$
\struck{\ill{}, fibré vectoriel} \add{\uncertain{aura}}
\[
  (57) \qquad \widetilde{B} = \check{\mathbb{V}}(\widetilde{\mathfrak{L}}) \cap
  \underbrace{\mathbb{V}(\mathcal{M})^{*}}_{\mathcal{L}}
  \qquad (B = \widetilde{B}/\mathbb{G}_m)
\]
\note{les lettres $B$, $\widetilde{B}$, $B_1$ sont écrites au-dessus des mots qu'elles nomment ; la seconde (lue $\widetilde{B}$) porte peut-être une barre plutôt qu'un tilde. $\mathfrak{L}$ transcrit sa lettre ronde pour l'algèbre de Lie d'un Borel ; son indice 1 dans « $\mathfrak{L}_1$ » est peu net. « en car.\ $\neq 2$ » est souligné de deux traits. Deux traits verticaux dans la marge gauche tiennent l'énoncé de la Proposition. Le « $(B = \widetilde{B}/\mathbb{G}_m$ » final est coupé par le bord de la feuille.}

\page{77}
Pour qu'un sous-fibré vectoriel \add{$\mathfrak{L}$} de $\widetilde{\mathfrak{g}}$ de
corang~1 soit de la forme $\underline{\mathrm{Lie}}(B)$, il f.\ et s.\ que le fibré
\add{$\mathfrak{L}^{\perp}$} \struck{\ill{}} de rang~1 orthogonal
\add{$\mathfrak{L}^{\perp}$} dans $\widetilde{\mathfrak{g}}$ satisfasse
\[
  (q \,|\, \mathfrak{L}^{\perp}) = 0
\]
(et ceci, bien sûr, caractérise $q$ \add{(à scalaire \uncertain{inversible} près)} --
\add{en termes de $\widetilde{G}$} -- \ill{} \ill{} \ill{} signe près \add{sur
$\widetilde{\mathfrak{g}}$ (\uncertain{définissant} $\check{\widetilde{\mathfrak{g}}}$)}
-- \ill{} \ill{} qu'il suffit de définir \ill{} \ill{} \uncertain{de façon}
compatible avec \ill{} \uncertain{de base})
\marginal{\emph{Question} : suffit-il pour $\mathfrak{L}$ \uncertain{isotrope} \ill{} de
\ill{} ss-\ill{} \add{$\mathfrak{L}$} i.e.\ stable par crochet ??}
\marginal{\emph{NB} Je \uncertain{reprends} \ill{} \ill{} dire que $\mathfrak{L}$ est
formé d'él.\ de \emph{\uncertain{carré nul}} de $\mathcal{M}$}
\note{les deux notes marginales sont écrites en oblique dans la marge gauche ; la première est séparée du texte par un trait vertical, la seconde est reliée au texte par un long trait. Les lettres $\mathfrak{L}$ et $\mathfrak{L}^{\perp}$ sont écrites au-dessus de la ligne, la seconde deux fois. La fin de la parenthèse, chargée d'insertions interlinéaires, se lit mal.}

\emph{NB} Plaçons-nous \emph{en car.~2}. Alors
\[
  (58) \qquad \underline{O}_S 1 \subset \widetilde{\mathfrak{g}}
\]
et l'\struck{application} morphisme
\[
  (59) \qquad X \longrightarrow \mathbb{P}(\widetilde{\mathfrak{g}})
\]
\uncertain{allant} des Borel de $\widetilde{\mathfrak{g}}$ vers leurs algèbres de
Lie, considérées comme sous-fibrés vect.\ de corang~1, est celle qui à la section de
$X$, identifiée à une \add{fibré en} droite \struck{dans} de
$\widetilde{\mathfrak{g}}$ satisfaisant $q \,|\, L = 0$, associe son orthogonal
dans $\widetilde{\mathfrak{g}}$ pour $\varphi$. \struck{L'image} \add{Comme}
l'orthogonal contient $\underline{O}_S \cdot 1$, \struck{et dans les}
(qui est le noyau de $\varphi$ i.e.\ de $f$), \struck{et} le morphisme
précédent se factorise par
\note{dans (59), un premier mot avant $\mathbb{P}$ est biffé et illisible ; sous le « NB » de « Plaçons-nous », un second « NB » est couvert de hachures, suivi d'une lettre isolée.}

\page{78}
$\mathbb{P}(\widetilde{\mathfrak{g}}/\underline{O}_S)$, \struck{\ill{}} i.e.
\[
  (60) \qquad X \longrightarrow \mathbb{P}(\widetilde{\mathfrak{g}}/\underline{O}_S)
\]
qui est lui-même un morphisme de fibrés en droites projectives.
Si \struck{le morphisme} \add{(59) donc (60)} était encore un monomorphisme, il
\struck{aurait} résulterait que (60) est un isomorphisme, donc $X$ serait « banal »
dès qu'on est en car.~2 -- or il n'est pas vrai qu'en car.~2 les fibrés en droites
projectives soient tous banals !
En fait, on trouve ici un iso.\ canonique
\[
  (61) \qquad \mathbb{P}(\widetilde{\mathfrak{g}}/\underline{O}_S) \simeq X^{(2)}
  \quad (\text{frobénisé})
\]
et (60) n'est autre que le morphisme canonique
\[
  X \longrightarrow X^{(2)}
\]
qui est un épimorphisme radiciel (et plat \add{fini loc.\ libre} de rang $2^d = 2$),
mais en effet nullement un monomorphisme \ldots
\note{« frobénisé » est relié à l'exposant $(2)$ par une flèche. La page s'arrête ici, au bas d'une moitié laissée blanche : c'est la dernière du dossier.}

\end{document}
